% ===========================================================================
\subsection{\file{new-design-reference/index.html}}
% ===========================================================================

\textbf{HTML, 395 lines.} The page itself: every word a visitor reads and the
structure holding those words. Delete every other file and this one still opens in
a browser and still says everything the project wants to say. It would simply look
plain.

\begin{provenance}{new-design-reference/index.html}
Hand-written. It arrived complete in the first commit, \code{1f7ed12} on
2026-07-28, at the path \file{index.html}, and moved to its present path in
\code{0dd839e} on 2026-07-29 with no change to its contents. No generator banner,
no framework scaffold, and no \file{package.json} that could have produced it.

Its icons are a different matter and worth naming precisely. Every icon is an
inline SVG whose path data matches the Heroicons outline set, a widely used free
library: the shapes, the \code{24 24} view box, the \code{stroke-linecap} and
\code{stroke-linejoin} attributes and the \code{1.5}-based coordinates are that
library's house conventions. \textbf{The file carries no attribution comment}, so
this is an inference from the path data, not a documented fact.

Its skeleton is shared with the sibling repository \code{ai-digitalise}: identical
head block, identical navigation shape, identical metrics band, identical booking
panel and footer. Which was written first is not recorded in either repository and
is not guessed here; \code{ai-digitalise} was committed a few hours earlier on the
same day. \textbf{Confidence: certain for hand-written, inferred for the icon
origin, and the sibling relationship is stated as a resemblance with a shared
skeleton rather than as a direction of copying.}
\end{provenance}

\subsubsection*{Why it exists}

Without this file there is no page. The stylesheet would style class names
belonging to nothing, and the JavaScript would search for elements that do not
exist and quietly do nothing, which is exactly how it is written to behave.

\subsubsection*{Line by line}

\paragraph{Lines 1 to 17: the head, which the visitor never sees.}

\begin{lstlisting}[language=HTML,firstnumber=1,caption={\file{index.html}, lines 1--17}]
<!DOCTYPE html>
<html lang="en">
<head>
<meta charset="UTF-8" />
<meta name="viewport" content="width=device-width, initial-scale=1.0" />
<title>Salman Adnan: AI Systems for Entrepreneurs</title>
<meta name="description" content="I build AI systems and digital products that grow businesses without adding headcount. 17 projects, 1.6M+ organic views." />
<meta property="og:title" content="Salman Adnan" />
<meta property="og:description" content="AI systems that grow your business without adding headcount." />
<meta property="og:type" content="website" />
<meta name="theme-color" content="#05070d" />

<link rel="preconnect" href="https://fonts.googleapis.com" />
<link rel="preconnect" href="https://fonts.gstatic.com" crossorigin />
<link href="https://fonts.googleapis.com/css2?family=Space+Grotesk:wght@500;600;700&family=Inter:wght@400;500;600;700&display=swap" rel="stylesheet" />
<link rel="stylesheet" href="/css/site.css" />
</head>
\end{lstlisting}

Line 1 is not a tag. \code{<!DOCTYPE html>} declares ``interpret this as modern
HTML''. Without it, browsers fall back to a compatibility mode imitating software
from the 1990s. It must come first and has no closing partner.

Line 2 opens the document and declares its language English. That attribute does
real work: screen readers use it to choose a pronunciation and search engines use
it to decide who to show the page to.

Line 3 opens the \code{<head>}. Everything to line 17 is information \emph{about}
the page rather than content \emph{of} it, and none of it appears on screen.

Line 4 declares the character encoding UTF-8, the scheme letting a file contain any
character in any language. It matters concretely here because the page contains a
middle dot, which would otherwise render as nonsense. It should come as early as
possible, since the browser guesses until told.

\begin{jargon}{Meta tag}
A tag inside \code{<head>} carrying a fact about the page rather than content for
it. It has no closing partner and displays nothing. Different meta tags are read by
completely different audiences: browsers, search engines and social networks each
look for their own.
\end{jargon}

Line 5 is the viewport declaration, and it is the single line making the page work
on a phone. \code{width=device-width} tells the browser to treat the page width as
the real device width, and \code{initial-scale=1.0} says do not zoom out to begin
with. Without it a phone pretends to be about 980 pixels wide and shrinks
everything, which is why old sites look like tiny unreadable newspapers. Every
media query in the stylesheet depends on this line.

Line 6 is the title: the browser tab text and the blue headline in a search result.
Note it, because section 1.2 used it to prove this page is not the live one.

Line 7 is the description search engines often show as grey text under the title.
It is worth reading closely, because it makes two specific claims (``17 projects,
1.6M+ organic views'') that also appear on the page itself in two other places.

Lines 8 to 10 are Open Graph tags, a convention invented by Facebook and now read
by nearly every chat app to build the preview card when somebody pastes a link.

\begin{honest}{The Open Graph set is missing the tag that matters most}
Present: \code{og:title}, \code{og:description}, \code{og:type}. Missing:
\code{og:image} and \code{og:url}. Without an image, a pasted link produces a
preview card with no picture, which is the difference between a link that gets
clicked and one that does not.

This is more pointed here than in the sibling repository. The live salmanadnan.com
carries a complete set including \code{og:image} pointing at
\code{/assets/og.png}, verified on 2026-08-11. So whoever built the live version
supplied exactly the asset this one lists as required and never produced. See
section 7.5.
\end{honest}

Line 11 sets \code{theme-color} to \code{\#05070d}, the same near-black as the page
background, so a mobile browser paints its address bar to match and the browser
chrome blends into the page instead of framing it in white.

\begin{honest}{The theme colour was not rebranded}
\code{\#05070d} is the shared near-black, which is correct. But note it is
identical to the sibling's, as is the entire head block apart from the title and
description. The rebrand described in commit \code{05210af} touched the stylesheet
and left this file's shared parts alone, which is consistent with the five missed
colour literals documented in section 7.2.
\end{honest}

Lines 13 and 14 are \code{preconnect} hints, telling the browser it will shortly
need these two Google servers so it should start the handshake now. Establishing a
secure connection takes several network round trips, and doing that early typically
saves a few hundred milliseconds. The \code{crossorigin} attribute on line 14 is
required: font files are fetched under stricter rules than stylesheets, and a
preconnect made without it opens a connection the font request cannot reuse. That
subtlety catches out experienced developers constantly.

Line 15 fetches the fonts: Space Grotesk in three weights and Inter in four.
\code{display=swap} at the end is the important part. It tells the browser to render
text immediately in a fallback typeface and swap in the real one when it arrives,
rather than hiding text for up to three seconds. Text reflows slightly; the benefit
is that a visitor on a slow connection can read the page at all.

\begin{jargon}{Font weight}
How heavy a typeface looks, on a numeric scale where 400 is normal and 700 is bold.
Requesting \code{wght@500;600;700} downloads three separate files, which is why the
list is kept short.
\end{jargon}

Line 16 links the stylesheet, with the leading slash discussed in section 11: correct
when served by a web server, wrong when the file is opened off a disk.

Line 17 closes the head.

\paragraph{Lines 18 to 28: the navigation bar.}

\begin{lstlisting}[language=HTML,firstnumber=18,caption={\file{index.html}, lines 18--28}]
<body>

<!-- ══════════════════════════════════════════════
     NAV
══════════════════════════════════════════════ -->
<nav class="nav" aria-label="Site navigation">
  <div class="nav-inner">
    <span class="nav-logo">Salman Adnan</span>
    <a href="#book" class="btn btn-primary">Book a call</a>
  </div>
</nav>
\end{lstlisting}

Line 18 opens \code{<body>}. Everything to line 394 is what a visitor sees.

Lines 20 to 22 are a comment banner. Anything between \code{<!-{}-} and \code{-{}->} is
ignored entirely by the browser. This file uses six of these as signposts, one per
band, each drawn with a rule of repeated double-line characters. They are the
author's table of contents and the reason a 395-line file stays navigable.

Line 23 opens the navigation. \code{<nav>} is a \emph{semantic} tag: it looks no
different from a plain \code{<div>} but announces to screen readers and search
engines that this is the site's navigation. \code{aria-label} gives that region a
spoken name.

\begin{jargon}{ARIA}
Accessible Rich Internet Applications: attributes beginning \code{aria-} that
describe a page to assistive software. They change nothing visually.
\code{aria-label} supplies a name; \code{aria-hidden="true"} tells assistive
software to skip something because it is purely decorative.
\end{jargon}

Line 24 opens a \code{<div>} labelled \code{nav-inner}. A \code{<div>} is a generic
container with no meaning, existing purely to be styled. The reason for a second
container is a standard layout pattern: the outer element spans the full window so
its background stretches edge to edge, while the inner one is capped and centred so
the logo and button line up with the content below.

\begin{honest}{The logo is not a link}
Line 25 is a \code{<span>} containing the words ``Salman Adnan''. It is not
clickable, so there is no way back to the top of the page by clicking it, which is
the behaviour every visitor expects from a site logo.

On a single-page site this is defensible, and it differs from the sibling
repository, where the logo is an \code{<a>} pointing at \code{/} with a drawn mark
beside it. That divergence is one of the few places where the two skeletons
genuinely part company, and this is the weaker of the two.
\end{honest}

Line 26 is the navigation's only button, a link to \code{\#book}. It is a link and
not a \code{<button>}, which is correct: it navigates rather than performing an
action. It carries two classes, \code{btn} for the shared shape and
\code{btn-primary} for the filled variant. That two-class pattern repeats
throughout and is explained in section 7.2.

Lines 27 and 28 close the containers.

\paragraph{Lines 29 to 46: the hero opens, and the placeholder that defines this
project's state.}

\begin{lstlisting}[language=HTML,firstnumber=29,caption={\file{index.html}, lines 29--46}]


<!-- ══════════════════════════════════════════════
     HERO
     UPDATE: replace .hero-photo-placeholder with:
       <img src="assets/salman.jpg" alt="Salman Adnan" style="width:100%;height:100%;object-fit:cover;object-position:top;">
══════════════════════════════════════════════ -->
<section class="hero" aria-label="Hero">

  <!-- photo side -->
  <div class="hero-photo-side">
    <div class="hero-photo-frame">

      <!-- PLACEHOLDER: swap for Salman's actual photo (see comment above) -->
      <div class="hero-photo-placeholder" aria-hidden="true">
        <div class="placeholder-initial">S</div>
        <div class="placeholder-note">Drop Salman's photo here<br>assets/salman.jpg</div>
      </div>
\end{lstlisting}

Lines 31 to 35 are the second banner, and it carries an instruction rather than
just a label: replace the placeholder with an \code{<img>} tag, given in full,
including the exact styling needed (\code{object-fit: cover} to fill the frame
without distortion and \code{object-position: top} to keep a face in shot when the
frame is cropped). Whoever wrote this knew precisely what was needed and did not
have the file.

Line 36 opens the hero. Line 39 opens the left column and line 40 the frame that
holds the photograph and the badges floating over it.

Lines 42 to 46 are the most consequential nine lines in the repository.

\begin{honest}{Half the opening screen is a placeholder addressed to the developer}
The left column of the hero, which on a desktop is half of everything visible when
the page loads, contains no photograph. It contains a giant letter \code{S} in an
orange gradient and, beneath it, the words \emph{``Drop Salman's photo here''}
followed by \emph{``assets/salman.jpg''}.

That is not a code comment. It is content: it sits inside the page, is styled by
\file{site.css} lines 173 to 179, and renders on screen. The \code{aria-hidden} on
line 43 hides it from screen readers, which is correct for a decorative gradient
and means a screen-reader user is spared the instruction, while every sighted
visitor reads it.

For a personal brand page, whose entire proposition is a person, the missing
element is the person. This single fact is why the honest description of this
project is ``unfinished'' rather than ``finished and parked''. The sibling
\code{ai-digitalise} has no equivalent: its placeholders are gradients that merely
look like unfinished thumbnails, not text asking for a file.
\end{honest}

Note that the note names the path \code{assets/salman.jpg} while the instruction in
the banner comment above names \code{assets/salman.jpg} and \file{ASSETS.md} names
\code{/assets/salman.jpg} with a leading slash. The leading slash is the one that
would actually work, for the same reason as the stylesheet path.

\paragraph{Lines 47 to 59: the first floating badge.}

\begin{lstlisting}[language=HTML,firstnumber=47,caption={\file{index.html}, lines 47--59}]

      <!-- floating result badges over the photo -->
      <div class="hero-float-badge" aria-hidden="true">
        <div class="badge-icon badge-icon-blue">
          <svg viewBox="0 0 24 24" fill="none" stroke="#3f82ff" stroke-width="2" aria-hidden="true">
            <path stroke-linecap="round" stroke-linejoin="round" d="M3 13.125C3 12.504 3.504 12 4.125 12h2.25c.621 0 1.125.504 1.125 1.125v6.75C7.5 20.496 6.996 21 6.375 21h-2.25A1.125 1.125 0 013 19.875v-6.75zM9.75 8.625c0-.621.504-1.125 1.125-1.125h2.25c.621 0 1.125.504 1.125 1.125v11.25c0 .621-.504 1.125-1.125 1.125h-2.25a1.125 1.125 0 01-1.125-1.125V8.625zM16.5 4.125c0-.621.504-1.125 1.125-1.125h2.25C20.496 3 21 3.504 21 4.125v15.75c0 .621-.504 1.125-1.125 1.125h-2.25a1.125 1.125 0 01-1.125-1.125V4.125z"/>
          </svg>
        </div>
        <div>
          <div class="badge-text-main">17 projects</div>
          <div class="badge-text-sub">delivered</div>
        </div>
      </div>
\end{lstlisting}

Three small frosted cards drift slowly over the photograph. Each has an icon in a
tinted square and two lines of text: a number and a caption.

Line 49 is the container, marked \code{aria-hidden} because the same claims appear
as real text further down the page. That is the correct call: repeating them to a
screen reader would be duplication.

Lines 51 to 53 draw a bar-chart icon. Note \code{stroke="\#3f82ff"}, a literal
blue, rather than a token.

\begin{honest}{The badge icons are the one place blue is deliberate}
Three badges use three different colours: blue here, orange on line 63, green on
line 75, each matching a tinted background defined at \file{site.css} lines 230 to
232. Used as a set of three distinct signal colours this is a reasonable design
choice, and it is the only appearance of blue in this file that is plainly
intentional rather than left over from the sibling.

It is still written as three literals where two tokens exist (\code{-{}-accent2} is
exactly this blue and is otherwise never used, and \code{-{}-green} is exactly that
green). Section 7.2 covers what that costs.
\end{honest}

Lines 55 to 58 hold the text: ``17 projects'' and ``delivered''. That number will
appear twice more before the page ends.

\paragraph{Lines 60 to 71: the second badge.}

\begin{lstlisting}[language=HTML,firstnumber=60,caption={\file{index.html}, lines 60--71}]

      <div class="hero-float-badge" aria-hidden="true">
        <div class="badge-icon badge-icon-orange">
          <svg viewBox="0 0 24 24" fill="none" stroke="#ff6719" stroke-width="2" aria-hidden="true">
            <path stroke-linecap="round" stroke-linejoin="round" d="M2.25 18L9 11.25l4.306 4.307a11.95 11.95 0 015.814-5.519l2.74-1.22m0 0l-5.94-2.28m5.94 2.28l-2.28 5.941"/>
          </svg>
        </div>
        <div>
          <div class="badge-text-main">1.6M+ views</div>
          <div class="badge-text-sub">organic, for clients</div>
        </div>
      </div>
\end{lstlisting}

Structurally identical, with an upward-trending chart icon in orange and the claim
``1.6M+ views / organic, for clients''. The orange \code{\#ff6719} here is the
brand accent written literally rather than as \code{var(-{}-accent)}.

\paragraph{Lines 72 to 86: the third badge, and the end of the photo column.}

\begin{lstlisting}[language=HTML,firstnumber=72,caption={\file{index.html}, lines 72--86}]

      <div class="hero-float-badge" aria-hidden="true">
        <div class="badge-icon badge-icon-green">
          <svg viewBox="0 0 24 24" fill="none" stroke="#34d399" stroke-width="2" aria-hidden="true">
            <path stroke-linecap="round" stroke-linejoin="round" d="M9 12.75L11.25 15 15 9.75M21 12a9 9 0 11-18 0 9 9 0 0118 0z"/>
          </svg>
        </div>
        <div>
          <div class="badge-text-main">300% avg growth</div>
          <div class="badge-text-sub">across client sites</div>
        </div>
      </div>

    </div>
  </div><!-- /photo side -->
\end{lstlisting}

The third badge uses the tick-in-a-circle icon and the claim ``300\% avg growth /
across client sites''. Its green is \code{\#34d399}, which is exactly the value of
the \code{-{}-green} token defined at \file{site.css} line 15 and written literally
here.

Line 85 closes the frame and line 86 closes the column, with a trailing comment
naming what was closed. That habit is worth copying: in a file with this much
nesting a bare \code{</div>} tells the reader nothing.

\begin{keyidea}{Why these three badges are numbered from 2 in the stylesheet}
The stylesheet positions the badges with \code{.hero-float-badge:nth-child(2)},
\code{(3)} and \code{(4)}, at \file{site.css} lines 217 to 219. They are children
2, 3 and 4 of the frame because the placeholder on line 43 is child 1.

This survives the intended fix: replace the placeholder with an \code{<img>} as
the comment instructs and the image becomes child 1, leaving the badges where they
are. That is either careful thinking or luck; either way the numbers do not need
to change when the photograph arrives, which is worth knowing before somebody
"corrects" them.
\end{keyidea}

\paragraph{Lines 87 to 112: the right column, down to the buttons.}

\begin{lstlisting}[language=HTML,firstnumber=87,caption={\file{index.html}, lines 87--112}]

  <!-- content side -->
  <div class="hero-content-side">
    <div class="hero-tag">
      <svg width="14" height="14" viewBox="0 0 24 24" fill="none" stroke="currentColor" stroke-width="2.5" aria-hidden="true">
        <path stroke-linecap="round" stroke-linejoin="round" d="M9.813 15.904L9 18.75l-.813-2.846a4.5 4.5 0 00-3.09-3.09L2.25 12l2.846-.813a4.5 4.5 0 003.09-3.09L9 5.25l.813 2.846a4.5 4.5 0 003.09 3.09L15.75 12l-2.846.813a4.5 4.5 0 00-3.09 3.09z"/>
      </svg>
      Digitalise Agency
    </div>

    <h1 class="hero-h1">
      Salman<br>
      <em>Adnan.</em>
    </h1>

    <p class="hero-descriptor">I build AI systems that grow<br>your business without adding headcount.</p>

    <div class="hero-actions">
      <a href="#book" class="btn btn-primary">
        Book a call
        <svg width="16" height="16" viewBox="0 0 24 24" fill="none" stroke="currentColor" stroke-width="2.5" aria-hidden="true">
          <path stroke-linecap="round" stroke-linejoin="round" d="M13.5 4.5L21 12m0 0l-7.5 7.5M21 12H3"/>
        </svg>
      </a>
      <a href="#work" class="btn btn-ghost">See the work</a>
    </div>
\end{lstlisting}

Line 89 opens the content column. Lines 90 to 95 are the small pill above the name,
containing a sparkle icon and the words ``Digitalise Agency'', which ties this
personal page back to the agency.

Lines 97 to 100 are the headline in an \code{<h1>}. A page should have exactly one,
and this page does: it is the document's title as far as a screen reader or a
search engine is concerned, and here it is simply the person's name.
\code{<br>} forces the line break so the surname lands on its own line.

\begin{honest}{\code{<em>} is used for colour, not for emphasis}
Line 99 wraps ``Adnan.'' in \code{<em>}, and \file{site.css} line 272 then removes
the italics that tag normally implies and colours the word orange instead.
\code{<em>} means emphasis and a screen reader may change its intonation for it.
Using it as a colour hook works and slightly misrepresents the content to assistive
software; a \code{<span>} with a class would have been the honest choice. The
sibling repository does exactly the same thing, so it is a habit rather than a
slip. The intent is plainly visual.
\end{honest}

Line 102 is the tagline, and it repeats the page description from line 7 almost
word for word. \code{<br>} again controls where it breaks.

Lines 104 to 112 hold the two buttons: the primary one to \code{\#book} containing
both text and an arrow icon, and the ghost one to \code{\#work}. Both are ordinary
links, and both are intercepted by \file{main.js} line 55 to scroll smoothly rather
than jump.

\paragraph{Lines 113 to 130: the three small statistics, and the end of the hero.}

\begin{lstlisting}[language=HTML,firstnumber=113,caption={\file{index.html}, lines 113--130}]

    <div class="hero-mini-stats">
      <div class="hero-mini-stat">
        <div class="hero-mini-num">17+</div>
        <div class="hero-mini-lbl">Projects delivered</div>
      </div>
      <div class="hero-mini-stat">
        <div class="hero-mini-num">1.6M+</div>
        <div class="hero-mini-lbl">Organic views</div>
      </div>
      <div class="hero-mini-stat">
        <div class="hero-mini-num">10K+</div>
        <div class="hero-mini-lbl">Users served</div>
      </div>
    </div>
  </div>

</section>
\end{lstlisting}

Three statistics under a hairline: 17+ projects, 1.6M+ organic views, 10K+ users
served.

\begin{honest}{The same three numbers exist twice on this page, in two different
formats}
These three are written as plain visible text. The metrics band at lines 258 to 279
states the same three (plus a fourth) as \code{data-target} attributes with a
visible \code{0}, to be counted up by JavaScript.

So updating a number means editing it in two places, in two different ways, and
nothing checks that they agree. \file{ASSETS.md} line 24 knows this and says so
explicitly, which means it was a recognised cost rather than an oversight.

There is an upside worth stating. Because the hero copies are plain text, they are
the only version a search engine or a visitor without JavaScript ever sees, so this
duplication accidentally rescues the page from the ``claims zero of everything''
failure that the sibling repository has in full. The sibling has no hero
statistics, so its numbers vanish entirely without JavaScript.
\end{honest}

Line 130 closes the hero section.

\paragraph{Lines 131 to 162: the work grid, and the first project card.}

\begin{lstlisting}[language=HTML,firstnumber=131,caption={\file{index.html}, lines 131--162}]


<!-- ══════════════════════════════════════════════
     WORK PORTFOLIO
     UPDATE: replace gradient .wt* classes with real thumbnails.
             Add poster="assets/thumb-N.jpg" to each tile img.
══════════════════════════════════════════════ -->
<section class="work-section" id="work" aria-label="Work portfolio">
  <div class="container">
    <div class="reveal" style="margin-bottom:0;">
      <p class="section-kicker">The work</p>
      <h2 class="section-title">What I've built<br>for clients.</h2>
    </div>

    <div class="work-grid stagger">

      <div class="work-card" style="--i:0">
        <div class="work-thumb wt1">
          <div class="work-badge">300% traffic</div>
          <div class="work-thumb-overlay">
            <div class="work-play">
              <svg viewBox="0 0 24 24" fill="white" aria-hidden="true">
                <path d="M5.25 5.653c0-.856.917-1.398 1.667-.986l11.54 6.348a1.125 1.125 0 010 1.971l-11.54 6.347a1.125 1.125 0 01-1.667-.985V5.653z"/>
              </svg>
            </div>
          </div>
        </div>
        <div class="work-info">
          <div class="work-client">Hi-Tec Bearings</div>
          <div class="work-title">5,000+ SKUs, 300% organic traffic growth, 45% faster page loads</div>
        </div>
      </div>
\end{lstlisting}

Lines 133 to 137 are the third banner, again carrying an instruction: replace the
gradient classes with real thumbnails.

Line 138 carries \code{id="work"}, which is what makes the hero's ``See the work''
button work. An \code{id} names one element, and a link to \code{\#work} scrolls to
whatever carries it.

\begin{jargon}{id, and how it differs from class}
An \code{id} names exactly one element and must be unique. A \code{class} labels
any number. Use an \code{id} when something needs to be linked to; use a
\code{class} when you want to style a group. This page has exactly two ids,
\code{work} and \code{book}, and both exist to be jumped to.
\end{jargon}

Lines 140 to 143 establish a pattern used by the remaining sections: a wrapper with
the class \code{reveal} holding a small uppercase kicker and a large heading. The
\code{reveal} class is what \file{main.js} line 50 watches for, fading the pair
upward when they scroll into view. The inline \code{margin-bottom:0} on line 140
overrides the 52-pixel margin the stylesheet gives \code{.section-title}, because
this section wants its grid closer than the others do.

Line 142 is an \code{<h2>}. Heading levels form the document's outline: the single
\code{<h1>} in the hero, then an \code{<h2>} per major section. Getting that order
right is how a screen-reader user navigates a long page, and this file gets it
right.

Line 145 gives the grid the class \code{stagger}, and each card carries
\code{style="-{}-i:0"} and so on. That custom property is multiplied by 90
milliseconds at \file{site.css} line 86 to delay each child's fade-in, so the cards
appear in sequence rather than together. Setting a custom property inline is the
cleanest way to pass a per-element number to CSS.

Each card is two halves. The top, \code{work-thumb}, is where a real thumbnail
would go and is currently a gradient from the class \code{wt1}. Over it sit a badge
and an overlay containing a play button, which the stylesheet keeps invisible until
the mouse is over the card. The bottom, \code{work-info}, holds the client name and
the result.

\begin{honest}{The play buttons promise videos that do not exist}
Every card shows a play button on hover and none plays anything: the card is not a
link and has no click handler.

This is milder than the same fault in the sibling repository, which also sets
\code{cursor: pointer} so the mouse changes to a hand. Here it does not, so the
promise is visual only. The cheapest honest fix in both is to remove the play
overlay until the videos exist.
\end{honest}

\paragraph{Lines 163 to 196: project cards two and three.}

\begin{lstlisting}[language=HTML,firstnumber=163,caption={\file{index.html}, lines 163--196}]

      <div class="work-card" style="--i:1">
        <div class="work-thumb wt2">
          <div class="work-badge">3x leads</div>
          <div class="work-thumb-overlay">
            <div class="work-play">
              <svg viewBox="0 0 24 24" fill="white" aria-hidden="true">
                <path d="M5.25 5.653c0-.856.917-1.398 1.667-.986l11.54 6.348a1.125 1.125 0 010 1.971l-11.54 6.347a1.125 1.125 0 01-1.667-.985V5.653z"/>
              </svg>
            </div>
          </div>
        </div>
        <div class="work-info">
          <div class="work-client">DigiMarvixx</div>
          <div class="work-title">3x WhatsApp leads per month, 60% mobile traffic capture</div>
        </div>
      </div>

      <div class="work-card" style="--i:2">
        <div class="work-thumb wt3">
          <div class="work-badge">1.2s load</div>
          <div class="work-thumb-overlay">
            <div class="work-play">
              <svg viewBox="0 0 24 24" fill="white" aria-hidden="true">
                <path d="M5.25 5.653c0-.856.917-1.398 1.667-.986l11.54 6.348a1.125 1.125 0 010 1.971l-11.54 6.347a1.125 1.125 0 01-1.667-.985V5.653z"/>
              </svg>
            </div>
          </div>
        </div>
        <div class="work-info">
          <div class="work-client">Synergy Event</div>
          <div class="work-title">800 participants, 1.2s load, 98/100 Lighthouse, 100% mobile score</div>
        </div>
      </div>
\end{lstlisting}

Structurally identical to the first, differing in the stagger index, the gradient
class, the badge text, the client and the result. The six clients across this grid
are Hi-Tec Bearings, DigiMarvixx, Synergy Event, IBA GPT, Prestige Sleep and IBA
CED, the same six as in the sibling repository, described here in slightly more
detail (this version adds ``45\% faster page loads'' and ``100\% mobile score'').

These are specific, named, checkable claims, which is a meaningfully different kind
of statement from the floating badges above them.

\paragraph{Lines 197 to 251: project cards four, five and six.}

\begin{lstlisting}[language=HTML,firstnumber=197,caption={\file{index.html}, lines 197--251}]

      <div class="work-card" style="--i:3">
        <div class="work-thumb wt4">
          <div class="work-badge">AI chat</div>
          <div class="work-thumb-overlay">
            <div class="work-play">
              <svg viewBox="0 0 24 24" fill="white" aria-hidden="true">
                <path d="M5.25 5.653c0-.856.917-1.398 1.667-.986l11.54 6.348a1.125 1.125 0 010 1.971l-11.54 6.347a1.125 1.125 0 01-1.667-.985V5.653z"/>
              </svg>
            </div>
          </div>
        </div>
        <div class="work-info">
          <div class="work-client">IBA GPT</div>
          <div class="work-title">AI university assistant answering student questions around the clock</div>
        </div>
      </div>

      <div class="work-card" style="--i:4">
        <div class="work-thumb wt5">
          <div class="work-badge">24/7 sales</div>
          <div class="work-thumb-overlay">
            <div class="work-play">
              <svg viewBox="0 0 24 24" fill="white" aria-hidden="true">
                <path d="M5.25 5.653c0-.856.917-1.398 1.667-.986l11.54 6.348a1.125 1.125 0 010 1.971l-11.54 6.347a1.125 1.125 0 01-1.667-.985V5.653z"/>
              </svg>
            </div>
          </div>
        </div>
        <div class="work-info">
          <div class="work-client">Prestige Sleep</div>
          <div class="work-title">Clinic booking and product shop combined, sells after hours</div>
        </div>
      </div>

      <div class="work-card" style="--i:5">
        <div class="work-thumb wt6">
          <div class="work-badge">4 dashboards</div>
          <div class="work-thumb-overlay">
            <div class="work-play">
              <svg viewBox="0 0 24 24" fill="white" aria-hidden="true">
                <path d="M5.25 5.653c0-.856.917-1.398 1.667-.986l11.54 6.348a1.125 1.125 0 010 1.971l-11.54 6.347a1.125 1.125 0 01-1.667-.985V5.653z"/>
              </svg>
            </div>
          </div>
        </div>
        <div class="work-info">
          <div class="work-client">IBA CED</div>
          <div class="work-title">Analytics platform with 4 executive dashboards tracking entrepreneurship outcomes</div>
        </div>
      </div>

    </div>
  </div>
</section>
\end{lstlisting}

The final three complete the set, with indices 3, 4 and 5 and gradients \code{wt4}
through \code{wt6}. Two of the six are for the same institution (IBA GPT and IBA
CED), the only repetition among them.

Note that the same play-triangle path data appears six times in this section and
nowhere is it shared. That is the cost of having no templating system, and it is
also the correct trade for a project that deliberately has no build step.

Line 249 closes the grid and lines 250 and 251 the container and section.

\paragraph{Lines 252 to 279: the metrics band.}

\begin{lstlisting}[language=HTML,firstnumber=252,caption={\file{index.html}, lines 252--279}]


<!-- ══════════════════════════════════════════════
     METRICS
     UPDATE: replace data-target values with real numbers
══════════════════════════════════════════════ -->
<section class="metrics-section" aria-label="Results">
  <div class="container">
    <div class="metrics-grid stagger">
      <div class="metric-item" style="--i:0">
        <div class="metric-num" data-target="17" data-suffix="+">0</div>
        <div class="metric-label">Projects Delivered</div>
      </div>
      <div class="metric-item" style="--i:1">
        <div class="metric-num" data-target="1.6" data-suffix="M+">0</div>
        <div class="metric-label">Organic Views for Clients</div>
      </div>
      <div class="metric-item" style="--i:2">
        <div class="metric-num" data-target="10" data-suffix="K+">0</div>
        <div class="metric-label">Users Served</div>
      </div>
      <div class="metric-item" style="--i:3">
        <div class="metric-num" data-target="300" data-suffix="%">0</div>
        <div class="metric-label">Avg. Traffic Growth</div>
      </div>
    </div>
  </div>
</section>
\end{lstlisting}

This band is where the JavaScript does its most visible work, and it introduces the
mechanism the counter depends on.

Look at line 262: the element's visible content is the single character \code{0},
and the real number lives in \code{data-target="17"}, with \code{data-suffix="+"}
beside it.

\begin{jargon}{Data attribute}
Any attribute whose name starts with \code{data-}. HTML ignores it completely and
displays nothing; it exists so JavaScript can attach its own information to an
element. \file{main.js} reads these two at lines 17 and 18 to know what to count
towards and what to print after it.
\end{jargon}

Why is the visible content \code{0}? Because that is what a visitor sees before the
counter runs, and it is also what they see forever if JavaScript is disabled or
fails. The robust alternative is to write the real number in the HTML and count
towards it from zero, which fails gracefully. This file chose the fragile form, and
is partly rescued by the hero statistics repeating three of the four as plain text.

This band still carries an \code{UPDATE} comment on line 256 telling the reader to
replace the values with real numbers, even though commit \code{692e162} claims to
have replaced all placeholder content with real project data.

\paragraph{Lines 280 to 318: the case studies, and the first card.}

\begin{lstlisting}[language=HTML,firstnumber=280,caption={\file{index.html}, lines 280--318}]


<!-- ══════════════════════════════════════════════
     CASE STUDIES (visual-only, numbers speak)
     UPDATE: replace gradient backgrounds with real screenshots.
             Use background-image on .case-visual or swap the div for <img>.
══════════════════════════════════════════════ -->
<section class="cases-section" aria-label="Case studies">
  <div class="container">
    <div class="reveal" style="margin-bottom:52px;">
      <p class="section-kicker">Case studies</p>
      <h2 class="section-title">Results in numbers.</h2>
    </div>

    <!-- Case 1 -->
    <div class="case-card reveal">
      <div class="case-visual cv1">
        <div class="case-visual-inner">
          <div class="case-visual-metrics">
            <div class="cvm">
              <span class="cvm-num">300%</span>
              <span class="cvm-lbl">Traffic growth</span>
            </div>
            <div class="cvm">
              <span class="cvm-num">5K+</span>
              <span class="cvm-lbl">Products</span>
            </div>
            <div class="cvm">
              <span class="cvm-num">99.9%</span>
              <span class="cvm-lbl">Uptime</span>
            </div>
          </div>
        </div>
      </div>
      <div class="case-content">
        <div class="case-client-name">Hi-Tec Bearings</div>
        <div class="case-title">Industrial e-commerce with 5,000+ SKUs that tripled organic traffic</div>
      </div>
    </div>
\end{lstlisting}

Two large cards, each a two-column arrangement of a visual panel and a block of
text. The visual is a gradient standing in for a screenshot, with three small
frosted chips of metrics laid over its bottom edge.

Line 295 gives the card the classes \code{case-card} and \code{reveal}, so it fades
in as a whole rather than being staggered.

Lines 299 to 310 are the three metric chips: 300\% traffic growth, 5K+ products,
99.9\% uptime. Note that 99.9\% uptime is a claim appearing nowhere else on the page
and in no other repository in this pair.

\paragraph{Lines 319 to 347: the second case study, and a dead class.}

\begin{lstlisting}[language=HTML,firstnumber=319,caption={\file{index.html}, lines 319--347}]

    <!-- Case 2 -->
    <div class="case-card case-card-reverse reveal">
      <div class="case-visual cv2" style="order:2">
        <div class="case-visual-inner">
          <div class="case-visual-metrics">
            <div class="cvm">
              <span class="cvm-num">3x</span>
              <span class="cvm-lbl">Leads/month</span>
            </div>
            <div class="cvm">
              <span class="cvm-num">60%</span>
              <span class="cvm-lbl">Mobile traffic</span>
            </div>
            <div class="cvm">
              <span class="cvm-num">WhatsApp</span>
              <span class="cvm-lbl">Integration</span>
            </div>
          </div>
        </div>
      </div>
      <div class="case-content" style="order:1">
        <div class="case-client-name">DigiMarvixx</div>
        <div class="case-title">Digital marketing site that tripled monthly leads through WhatsApp</div>
      </div>
    </div>

  </div>
</section>
\end{lstlisting}

The second card is meant to mirror the first: text on the left, visual on the
right. Line 321 gives it the extra class \code{case-card-reverse}, and lines 322
and 340 set \code{order:2} and \code{order:1} inline.

\begin{honest}{The reversal class matches nothing, and the layout works anyway}
\file{site.css} lines 436 and 437 define \code{.case-card.reverse}, which means ``an
element carrying both the class \code{case-card} and the class \code{reverse}''.
This element carries \code{case-card} and \code{case-card-reverse}. A hyphen does
not create a second class, so \textbf{those rules have never once applied}, and
neither has the third at line 559 in the responsive block.

The reversal happens regardless, because of the inline \code{order} properties on
lines 322 and 340, which tell the grid to place the visual second and the text
first. So the design is correct, three stylesheet rules are dead, and nothing
anywhere reported it.

This is the clearest example in the pair of what ``silent failure'' means in
practice: a whole mechanism can be wrong while the page looks right, because a
second mechanism quietly does the same job.
\end{honest}

\begin{jargon}{order}
A CSS property setting an item's position inside a flex or grid container
independently of where it sits in the HTML. \code{order:2} places an item after
one with \code{order:1}, regardless of which is written first.
\end{jargon}

Note also that the inline \code{style="order:2"} on line 322 sits alongside the
class \code{cv2}, so one element is being positioned by an inline style and
coloured by a class, which is the same mixed habit seen throughout the pair.

\paragraph{Lines 348 to 377: the booking panel.}

\begin{lstlisting}[language=HTML,firstnumber=348,caption={\file{index.html}, lines 348--377}]


<!-- ══════════════════════════════════════════════
     BOOK A CALL
     UPDATE: replace YOUR_CAL_LINK with your Cal.com or Calendly URL.
══════════════════════════════════════════════ -->
<section class="book-section" id="book" aria-label="Book a call">
  <div class="container">
    <div class="book-wrap reveal">
      <p class="book-eyebrow">Let's talk</p>
      <h2 class="book-title">30 minutes.<br>No slides.</h2>
      <p class="book-sub">Tell me about your content goals. I'll tell you if I can help.</p>
      <div class="book-cta-row">
        <!-- UPDATE: replace href with your booking link -->
        <a href="https://cal.com/salman-adnan/meeting" target="_blank" rel="noopener" class="btn btn-primary" style="padding:16px 36px;font-size:17px;">
          Book a call
          <svg width="18" height="18" viewBox="0 0 24 24" fill="none" stroke="currentColor" stroke-width="2.5" aria-hidden="true">
            <path stroke-linecap="round" stroke-linejoin="round" d="M13.5 4.5L21 12m0 0l-7.5 7.5M21 12H3"/>
          </svg>
        </a>
      </div>
      <p class="book-detail">
        <svg viewBox="0 0 24 24" fill="none" stroke="currentColor" stroke-width="2" aria-hidden="true">
          <path stroke-linecap="round" stroke-linejoin="round" d="M12 6v6h4.5m4.5 0a9 9 0 11-18 0 9 9 0 0118 0z"/>
        </svg>
        Free · <span>30 minutes · Zoom</span>
      </p>
    </div>
  </div>
</section>
\end{lstlisting}

This is the destination of every ``Book a call'' button, via \code{id="book"} on
line 354. It is the point of the whole document.

Line 356 combines two classes, \code{book-wrap} for the panel's appearance and
\code{reveal} so it fades in, the only place in the file where a layout class and a
behaviour class share an element.

Line 362 is the booking link, and three of its attributes matter.
\code{target="\_blank"} opens it in a new tab so the visitor does not lose the page.
\code{rel="noopener"} is a security measure that must accompany it: without it the
newly opened page receives a reference back to this one and can navigate it
somewhere else, a trick known as tab-nabbing. Its presence is correct and is the
kind of detail usually forgotten.

\begin{keyidea}{This booking link is real, and it outlived the design}
\code{cal.com/salman-adnan/meeting} is a genuine personal calendar, set
deliberately in commit \code{014406e}. The sibling repository's equivalent points
at \code{cal.com/davonex/30min}, another project's calendar, and is a known
outstanding placeholder.

More than that, this link is the one part of this design that is demonstrably still
in service: fetching the live salmanadnan.com on 2026-08-11 found the same link on
the page that replaced this one.
\end{keyidea}

\begin{honest}{The comment above the link is stale}
Line 352 still tells the reader to search for \code{YOUR\_CAL\_LINK}, and line 361
repeats the instruction. That literal placeholder was replaced on 2026-07-28 and
appears nowhere in the file. \file{ASSETS.md} line 28 repeats the same stale
instruction a third time. Somebody following the written instructions would search,
find nothing, and have no way to tell whether the job was done or the instruction
was wrong.
\end{honest}

Lines 369 to 374 are the reassurance line under the button: a clock icon, then
``Free'', then the duration and platform. The middle dots are the U+00B7 character
discussed in section 1.4 and render exactly.

\paragraph{Lines 378 to 395: the footer, and the script tag.}

\begin{lstlisting}[language=HTML,firstnumber=378,caption={\file{index.html}, lines 378--395}]


<!-- ══════════════════════════════════════════════
     FOOTER
══════════════════════════════════════════════ -->
<footer>
  <div class="container">
    <p>
      <a href="mailto:salman@digitalise.agency">salman@digitalise.agency</a>
      &nbsp;·&nbsp;
      <a href="https://digitalise.agency" target="_blank" rel="noopener">digitalise.agency</a>
    </p>
  </div>
</footer>

<script src="/js/main.js"></script>
</body>
</html>
\end{lstlisting}

The footer is one line: an email address and a link to the agency, separated by a
middle dot padded with \code{\&nbsp;} on each side.

\begin{jargon}{Non-breaking space}
\code{\&nbsp;} is a space a browser will not use as a place to wrap a line, and
which is not collapsed when several sit together. Here it does the second job:
ordinary spaces around the dot would collapse to one, and these keep the separator
visually padded.
\end{jargon}

Line 386 is a \code{mailto:} link, which opens the visitor's email program with the
address filled in. Line 388 links to the agency site, correctly carrying both
\code{target="\_blank"} and \code{rel="noopener"}.

This footer is better than the sibling's, which links to \code{digitalise.agency}
as text but points the link at \code{/}, its own home page. Here the text and the
destination agree.

Line 393 loads the JavaScript, and \textbf{where it sits is the whole point}. It is
the last thing before \code{</body>}, so by the time the browser reaches it every
element already exists. \file{main.js} begins searching for elements immediately,
with no wrapper telling it to wait, and that only works because of this placement.
Move this line into the \code{<head>} and the script would run against an empty
document, find nothing, and the page would quietly lose all three behaviours with
no error message.

Lines 394 and 395 close the body and the document.

\subsubsection*{What it connects to}

\textbf{This file names:} \file{new-design-reference/css/site.css} (line 16),
\file{new-design-reference/js/main.js} (line 393), Google Fonts (lines 13 to 15),
cal.com (line 362), and the agency site (line 388).

\textbf{This file is named by:} nothing inside the repository. The workflow copies
the whole folder rather than referring to files individually, and
\file{netlify.toml} publishes the folder as a whole. \file{ASSETS.md} discusses it
at length and cannot affect it.

\textbf{Internal links:} the navigation and hero both point at \code{\#book} (lines
26 and 105), and the hero also points at \code{\#work} (line 111). Those two ids,
on lines 354 and 138, are the only link targets in the file.

\subsubsection*{How it breaks}

\begin{itemize}
\item \textbf{Rename a class here without renaming it in the stylesheet} and that
  element loses its appearance silently. Nothing warns you. Section 7.2 documents a
  live instance of the reverse: a stylesheet class that matches nothing here.
\item \textbf{Move the \code{<script>} tag into the \code{<head>}} and all three
  behaviours stop, with no error. The metrics band would then permanently show four
  zeros.
\item \textbf{Insert an element before the placeholder on line 43} and all three
  floating badges jump, because the stylesheet positions them by child number.
\item \textbf{Open the file directly off a disk} and it renders unstyled, because of
  the leading slash on line 16. This looks like catastrophic breakage and is not.
\item \textbf{Disable JavaScript} and the page is fully readable but static. The
  metrics band shows four zeros; the three hero statistics still show their real
  values, which is what saves it.
\item \textbf{Delete the id on line 354} and every ``Book a call'' button silently
  does nothing, because \file{main.js} returns early when it cannot find a link's
  target.
\end{itemize}

\subsubsection*{Honest findings for this file}

The visible photo placeholder on lines 43 to 46; the missing \code{og:image} and
favicon; seven stale comments (lines 33, 42, 135, 256, 284, 352, 361); a
non-clickable logo on line 25; \code{<em>} used as a colour hook on line 99; the
same three numbers stated twice in two formats; play buttons that play nothing;
colour literals used where tokens exist (lines 51, 63, 75); and a
\code{case-card-reverse} class that no stylesheet rule matches.
